\documentclass[../main.tex]{subfiles}
\begin{document}
%==========================================TIÊU ĐỀ TẬP========================================
\begin{table}[h]
    \begin{adjustwidth}{-1cm}{}
        \begin{tabular}{>{\centering\arraybackslash}p{6cm}|>{\raggedright\arraybackslash}p{12.5cm}}
            \multirow{5}{6cm}
            % Cột bên trái: ÉPISODE và số tập
            {\\[-40pt]
                \flaregothic\fontsize{20pt}{20pt}\selectfont \centering
                \color{eptype}PERSONNALISÉ\color{black}\\[12pt]
                \vnmdisp\fontsize{36pt}{36pt}\selectfont
                \color{epnum}BỐN\color{black}\\[6pt]
                \cafeteria\fontsize{20pt}{20pt}\selectfont\color{epnum}(C-4)\color{black}
                \\[18pt]
                % \flaregothic\fontsize{14pt}{14pt}\selectfont
                % \color{epattr}BIENVENUE, \uline{\color{Banpire} Mel} !
                % \flaregothic\fontsize{18pt}{18pt}\selectfont
                % \color{epattr}ÉPISODE PERDU
                \color{black}
            }
             &
            % Dòng thứ nhất: tên tập tiếng Nhật
            {\fotskip\fontsize{18pt}{18pt}\selectfont るしあの60\ruby{秒}{びょ}\ruby{作}{さく}\ruby{戦}{せん}
            }                                                                                   \\[6pt]
            % Dòng thứ hai: tên tập tiếng Anh
             & {\cafeteria\fontsize{18pt}{18pt}\selectfont Baseball Madness}                                                                                                      \\[4pt]
            % Dòng thứ ba: tên tập tiếng Việt
             & {\vnmsans\fontsize{18pt}{18pt}\selectfont Cuộc tập dượt bóng chày}                                                                                                 \\[8pt]
            % Dòng thứ tư: Nhân vật xuất hiện
             & {\caviar{\fontsize{14pt}{14pt}\selectfont Nhân vật: \emp{kuon1} \emp{ryuuhou1} \emp{korone} \emp{pekora} \emp{rushia}}}
            \\[18pt]
            % % Dòng cuối cùng: Thời gian diễn ra của tập (thường sẽ là ngày phát hành)
            % & {\continuum\fontsize{16pt}{16pt}\selectfont Ngày phát hành: 12/09/2021}\\[18pt]
        \end{tabular}
    \end{adjustwidth}
\end{table}
%=========================================TRUYỆN CHÍNH========================================
\color{black}

\txtbx[2]
{{\color{vol} \uline{Tóm tắt:}} Rushia, Pekora và Korone tổ chức một trận bóng chày ba người tại văn phòng, với Kuon làm trọng tài bất đắc dĩ và Ryuuhou làm khán giả duy nhất. Pekora đánh trượt hai lượt, đến lượt thứ ba đánh trúng nhưng bị Korone bắt bóng, dẫn tới việc bị loại. Rushia đánh trúng và ghi điểm. Korone đánh bóng quá mạnh, làm vỡ cửa sổ, cả ba đuổi theo quả bóng. Kuon thở dài trước thiệt hại, Ryuuhou nhận xét kính cường lực cần sửa chữa. Tập này là chuyển thể từ thước phim "60 giây" (ÉPISODE 127), dựa trên các khung hình do Kanata, Lamy và Watame dựng trong đêm hỗn loạn.}

Giữa tháng Mười, hơi se lạnh của mùa thu đã len lỏi khắp nơi, những cơn gió mỏng manh thổi vào bên trong văn phòng \hll. Mặt trời dần lặn về phía chân trời phía tây, để lại những tia nắng xiên nhẹ trải dài trên mặt thảm, mang theo làn không khí mát lạnh của những ngày cuối thu.

Rushia đứng giữa căn phòng, tay cầm một quả bóng chày đã khá cũ. Cô mặc chiếc áo khoác mỏng màu xám tro, mái tóc xanh lơ được buộc cao như thường lệ, khẽ bay theo nhịp gió. Dáng vẻ nhỏ nhắn của cô trông chẳng hề phù hợp với môn thể thao mạnh mẽ này chút nào, thế nhưng ánh mắt lại toát lên sự nghiêm túc đến lạ thường.

``Bởi vì\dots'' cô thì thầm, giọng nhỏ nhẹ như đang tự nói với mình, ``nếu mình không chủ động làm gì đó, mọi người sẽ dần lãng quên mình mất\dots''

Không một ai nghe thấy câu nói vừa rồi. Nhưng chỉ vài phút sau, Pekora xuất hiện ở cửa phòng, tay cầm cây gậy nhựa màu hồng, vừa đi vừa nhai chiếc cà rốt trên tay. Bộ đồ thể thao của cô có hơi nhàu nhĩ, rõ ràng là chỉ mặc qua loa cho có bối cảnh, chứ không phải vì đam mê thể thao.

``Bà có chắc là muốn chơi thật đấy không, peko?'' cô hỏi, để gậy tựa trên vai, mắt nheo lại nhìn nàng Pháp sư, dường như đang nghi ngờ liệu có ai đó lại rủ mình chơi đánh cờ giữa trời mưa hay không.

Nàng Pháp sư gật đầu: ``Ừm. Chỉ cần chơi một trận duy nhất thôi-nanodesu.''

``Không có luật lệ kỳ lạ gì chứ, peko? Không được dùng bom, không được triệu hồi hồn ma, cũng không được kêu fan đến cổ vũ, đúng chứ?''

``Ừm,'' nàng Pháp sư lại gật đầu. ``Chỉ là một trận bóng chày bình thường thôi mà.''

Cùng lúc đó, từ cầu thang bên hông tòa nhà, Korone bước ra, vẫn đeo tai nghe và đeo chiếc ba lô trên vai. Cô ngáp dài, vươn vai thoải mái: ``Nghe nói em vận động nhẹ để giãn gân cốt à? Cho chị tham gia với!''

``Thế là đủ ba người rồi!'' Rushia thốt lên, vẻ mặt có chút vui mừng vì kế hoạch nhỏ của mình đang dần thành hiện thực.

``Khoan đã,'' Pekora ngẩng đầu lên, ``vậy ai làm trọng tài bây giờ, peko?''

Câu hỏi đó như một mũi tên xuyên qua cửa kính phòng kỹ thuật, nơi Kuon đang vừa nhâm nhi cà phê vừa kiểm tra bản dựng ứng dụng mới nhất, bên cạnh là Ryuuhou đang thong thả lướt điện thoại. Và dĩ nhiên, cậu Tổng bị lôi ra khỏi phòng ngay lập tức. ``Tự dưng gọi anh ra làm gì thế? Anh chỉ là một lập trình viên thôi, mấy cái thể thao hay vận động tay chân này anh không rành đâu!''

Nàng Trưởng nhóm bước ra sau anh trai, tay cầm ly trà ấm, nhìn người anh mình bị hai cô gái kéo lê với ánh mắt vô cùng bình thản, khẽ xắt xéo ông anh của mình: ``Để Onii làm trọng tài, khả năng cao là trận đấu chưa kết thúc mà trọng tài đã hụt hơi trước cả người chơi luôn.''

``Thế thì càng tốt!'' nàng Khuyển đẩy cậu ra giữa sân, nhét vào tay cậu một chiếc còi nhựa. ``Trọng tài tốt nhất chính là người không biết gì về môn thể thao này. Như vậy mới công bằng!''

``Đúng vậy đó peko!'' nàng Thỏ lập tức hưởng ứng, ``Chứ để Rushia làm trọng tài thì quả bóng còn có thể phát nổ nữa.''

``Không có đâu mà,'' nàng Pháp sư lẩm bẩm nhỏ.

Cậu Tổng đành thở dài, lắc đầu chấp nhận với đề xuất đến bất ngờ này. Trong lòng, cậu tự nhủ: ``\textit{Hy vọng trận cầu này sẽ không có gì vượt quá tầm kiểm soát\dots}''

Ryuuhou chọn một góc ghế xếp dưới bóng râm bên lề sân, nhấp nhẹ ngụm trà, tự nhận vai trò là khán giả duy nhất cho trận đấu sắp diễn ra. Một trận bóng chày ba người với cách phân vai chẳng mấy hợp lý cũng được sắp xếp xong. Rushia làm người ném bóng, Pekora làm người đánh, Korone làm người đuổi bắt bóng, và Kuon vô tình trở thành trọng tài trong khi chẳng mặn mà gì.

Một buổi chiều tháng Mười, chẳng ai biết lý do vì sao mọi người lại rảnh rỗi đến mức làm chuyện này. Có lẽ vì thời tiết hôm ấy quá đẹp, hoặc vì tâm trạng của ai đó cần một chút thư giãn. Hay có lẽ chỉ đơn giản là\dots\ ngày hôm đó, không có lịch ghi hình nào quan trọng cả.

\ct

``Vào vị trí hết đi!'' Kuon than thở, cậu chẳng khác gì một kẻ đang bị số phận bức bách, tay vẫn cầm chặt chiếc còi nhựa xanh lơ đã mẻ một góc. Cậu đứng rụt rè ở mép sân, luôn giữ một khoảng cách an toàn với ba cô gái kia, chẳng những sợ bị bóng chày đập trúng mặt, mà còn vì bên cạnh cậu là một nàng chó với khí chất như dao găm, bên kia là một nàng thỏ lúc nào cũng kêu ``peko peko'', và đứng giữa họ là một cô đồng luôn mang nỗi buồn vô tận. Làm sao cậu Tổng có thể tin rằng đây chỉ là một trận bóng chày bình thường được chứ?

Ở tâm sân, Rushia nắm chặt quả bóng trong lòng bàn tay. Cô hít một hơi thật sâu, nhắm chặt hai mắt như thể đang tập trung tụ linh khí từ cõi âm. Gió thổi nhẹ nhàng bay tà tóc cô, khiến Kuon không nhịn được mà nghĩ thầm: ``\textit{Chắc đây là lần đầu tiên trong đời mình thấy có người ném bóng mà trông giống đang triệu hồi hồn ma thế này.}''

Ryuuhou ngồi trên ghế xếp, khẽ nhõng mày lên và bình thản nhận xét: ``Cách ném bóng này trông giống nghi thức yểm bùa hơn là môn thể thao đấy, Rushia-san.''

``\textbf{Được rồi, bắt đầu lượt thứ nhất!}'' Kuon hô lớn, rồi thổi một tiếng còi vang lên. Rushia lấy đà, rồi tung người ném bóng. Cú ném không quá mạnh, nhưng quỹ đạo của quả bóng lại kỳ lạ vô cùng: nó xoắn một vòng nhẹ rồi lướt ngang trên mặt đất, cứ như đang trôi nổi lơ lửng giữa không trung.

``Ể\dots\ Ể!?'' Cả Pekora và Korone đều trố mắt ngạc nhiên trước đường đi trái với mọi vật lý của quả bóng.

Nàng Thỏ cố gắng vung mạnh cây chày về phía quả bóng ngay khi nó bay tới, nhưng\dots

*VÚT!*

``Peko!''

\dots cô đập hụt hoàn toàn, thậm chí còn bị văng ngã, mặt đập sầm xuống sàn nhà.

``\textbf{Nó đâu có đi thẳng đâu, peko!! Bà ném kiểu gì vậy!?}'' cô gượng dậy, vừa phủi bụi trên người vừa kêu lên.

``Làm sao mà tôi biết được!'' nàng Pháp sư chớp mắt liên hồi, bản thân cô cũng không ngờ mình vừa ném ra một cú bóng kỳ quái như vậy. ``Chắc là do gió thổi nó lệch đi thôi\dots''

Từ phía xa, Korone lao tới như một cơn gió, lăn một vòng trên đất và tóm gọn quả bóng trong tay, dù nó chỉ vừa mới chạm đất nhẹ nhàng. ``\textbf{Gâu gâu! Đuổi bóng là việc của chị mà\~{}}'' cô hét lớn, khiến bụi đất bay tung tóe khắp nơi.

Cậu Tổng lại thổi còi một lần nữa, giơ tay ra để thông báo: ``Lượt đầu tiên không tính! Bóng chưa qua được strike zone!''

``Strike zone là cái quái gì thế?'' cả ba cô gái đồng thanh hỏi lên.

``Ờ thì\dots\ là vùng mà người đánh có thể chạm được bóng ấy. Tóm lại là em đánh trượt rồi!'' Kuon gãi đầu gãi tai, cậu chả hiểu gì về các quy tắc của bóng chày cả.

``Thì cứ bảo là trượt đi là được mà!'' nàng Thỏ bực bội nói. ``Lại cứ nói cái từ sờ-trai-dôn kia làm gì\dots''

\ct

Đến lượt đánh thứ hai. Lần này, Rushia ném bóng nhanh hơn nhiều. Quả bóng bay vèo qua một đường parabol hoàn hảo, chẳng khác gì các cú ném chuyên nghiệp.

``Để em chuẩn bị tư thế nào\dots'' nàng Thỏ ngước nhìn quả bóng đang bay tới, vừa căn chỉnh lại vị trí của cây chày trong tay. Nhưng trước khi kịp vào tư thế, cô đã lùi nhanh hai bước như thể vừa thấy ma: ``\textbf{Peko!!}''

Quả bóng lúc này đã nằm gọn trong tay Korone. ``Sao thế, Peko-chan?''

``N-Nó nhìn em!!'' Pekora chỉ vào vệt bóng tối dưới chân mình. ``Em thề là cái bóng đó đang nhìn em, peko!!''

``Cái gì nhìn em cơ?'' Kuon nhướng mày thắc mắc.

``Cái bóng! Của quả bóng đó ấy, peko!!''

Ryuuhou dựa lưng vào thành ghế, điềm tĩnh bổ sung: ``Trong tâm lý học, người ta gọi đó là pareidolia, là hiện tượng não bộ tự tạo ra ảo giác về khuôn mặt từ những hình ảnh ngẫu nhiên đấy, Pekora-san. Chị có ngủ đủ giấc không vậy, hay là mắt chị bị nhìn gà hóa cuốc thế?''

``Đó là do ánh mặt trời đấy, đồ ngốc!'' Korone vừa nói vừa nhai nhóp nhép ít cát trên đất, trông như đang tập vai cho một bộ phim về chó nhẫn giả.

``Quả này cũng tính là trượt nhé,'' cậu Tổng tuyên bố, lần này giọng dứt khoát hơn hẳn.

\ct

Lượt đánh thứ ba, cũng là lượt cuối cùng của Pekora.

Nàng Pháp sư ném ra quả bóng quyết định. Cô lấy đà từ từ, chậm rãi như những cảnh phim thể thao học đường. Ngay khi tay vừa buông bóng, một bản nhạc nền tưởng tượng dường như vang lên trong đầu cô: một bản nhạc bi tráng mà chỉ có mình cô mới nghe thấy.

Bóng bay thẳng, mạnh mẽ và mượt mà trên không trung.

Pekora hét lên một tiếng xung kích ``\textbf{Peko!!!}'' rồi vung cây chày thật mạnh về phía quả bóng.

*CHÁT!*

Tiếng va chạm vang dội đến mức cậu trọng tài giật mình, làm rơi cả chiếc còi trên tay xuống đất.

Quả bóng bay lên cao, xoáy nhẹ giữa bầu trời. Cảnh tượng này dường như chuyển động chậm lại trong tâm trí của tất cả mọi người, ít nhất là theo lời kể của Rushia sau này.

``\textbf{Bóng bay rồi, peko!!}'' Pekora hô to.

``\textbf{Bóng bay kìa, gâu gâu!!}'' Korone lập tức lao theo, hai chân chạy thoăn thoắt, miệng vẫn không ngừng sủa nhỏ.

Kuon quay lưng lại, giơ tay ôm lấy đầu: ``Đừng để bóng đập vào mặt anh là được rồi\dots''

Sau một khoảnh khắc kéo dài như cả thế kỷ được thu nhỏ trong ba giây, Korone bật người cao lên, vươn tay ra với lấy quả bóng.

*PẶP!*

Nàng Khuyển bắt được bóng thành công.

``Yaaay!! Bắt được rồi, bắt được rồi!!'' cô nhảy tưng tưng vui vẻ, chẳng khác gì một chú chó săn vừa được thưởng món bánh quy yêu thích. Nàng Pháp sư mỉm cười hạnh phúc. Còn nàng Thỏ thì đứng sững người, cây chày vẫn còn giơ lên trên không. ``Peko\dots\ Em vừa đánh trúng nó rồi\dots\''

``Công nhận,'' Kuon vừa nhặt lại chiếc còi, vừa nói với niềm vui lây, nhưng rồi giọng cậu bỗng nghiêm túc hẳn, ``nhưng nếu đây là một trận bóng chày thật sự, thì em bị loại rồi đấy.''

Pekora ngoảnh đầu lại nhìn cậu, trố mắt ngạc nhiên: ``Ể? Nhưng mà em đã đánh trúng nó mà?''

``Đúng là đánh trúng, nhưng người ta đã bắt được bóng rồi,'' cậu Tổng cười khúc khích. ``Sau ba lần đánh trượt, em chính thức bị loại khỏi vòng đấu.''

Nàng Thỏ im lặng, như không thể tin vào những gì vừa nghe thấy. ``C-Cái gì chứ, peko\dots?''

``Tiếc quá nhỉ, Pekora-chan,'' Rushia cười trừ. ``Lần sau bà đánh mạnh hơn một chút là được thôi\dots''

Nàng Thỏ cúi đầu, rồi gào lên: ``\textbf{Tại sao lại thế chứ, peko!!!???}'' Cô nghiến chặt răng, ngẩng đầu lên nhìn hai người kia đang bụm miệng cười. Hai má cô phồng lên vì tức giận, vứt mạnh cây chày xuống sàn: ``Khỉ thật! Hai lượt đầu thì hụt, lượt cuối thì lại bị đối phương bắt bóng!''

``Nói chung thì trình của chị cũng chưa đủ để làm cầu thủ bóng chày được đâu, Pekora-san,'' Ryuuhou gọi vọng ra từ ghế, nét mặt không giấu được sự nham hiểm.

Rushia vừa lo lắng vừa thấy buồn cười: ``Nhưng Peko-chan cũng đã đánh trúng nó rồi mà\dots''

Korone vẫn đeo găng tay bắt bóng, híp mắt cười: ``Hehe\~{} nhưng mà bị bắt thì vẫn là bị loại đó nha\~{}''

``Lần sau em sẽ đánh cho bay luôn khỏi cái văn phòng này cho xem!'' nàng Thỏ tuyên bố đầy thách thức, vừa nói vừa lườm cả hai người còn lại.

Cậu Tổng (kiêm trọng tài) lật sổ ghi chép lại: ``Lượt thứ ba, Peko-chan đánh trúng. Koro-chan bắt được bóng. Như vậy Ru-chan vẫn chưa được tính điểm.''

Cậu khẽ liếc nhìn đồng hồ trên tay, rồi quay sang ba cô gái: ``Được rồi, đến lượt Ru-chan đánh rồi đấy. Ai làm pitcher - tức là người ném bóng - đây?''

Korone lập tức vẫy tay: ``Để chị ném cho!''

\ct

Rushia bước vào vị trí, cầm cây gậy lóng ngóng chẳng khác nào chưa từng đụng đến môn thể thao này bao giờ. Ở phía đối diện, Korone vẫn tươi cười toe toét, tay vung quả bóng chuẩn bị tung ra. Còn Pekora thì lóc cóc đứng đằng sau nàng Khuyển, gương mặt vẫn còn ngập tràn sự hậm hực sau lượt đánh vừa rồi.

``Á! Lại trượt nữa rồi!'' nàng Pháp sư kêu lên, còn quay vòng vòng tại chỗ vì tiếc nuối.

Korone ngửa cổ cười lớn: ``Lại thêm một lần nữa nha\~{}''

Đến lượt thứ hai, Rushia dồn toàn bộ sức lực, gồng mình lên và\dots

*CHÁT!*

Tiếng gậy va vào bóng vang lên rõ rệt. Quả bóng bay lên không trung nhưng lại xoáy ngang một cách kỳ lạ, đập mạnh vào bức tường rồi rơi thẳng xuống ngay trước chân Kuon.

Cậu Tổng cúi xuống nhặt quả bóng lên, thổi còi để thông báo: ``An toàn! Ru-chan được cộng điểm!''

Nàng Pháp sư thở dài một hơi, nhưng gương mặt nở ra nụ cười rạng rỡ: ``Công nhận chơi bóng chày thú vị thật đó-nanodesu!''

``Tôi thì chẳng thấy vui chút nào đâu, peko,'' đôi tai thỏ của Pekora cụp hẳn xuống, càng tức giận hơn khi thấy cô bạn cùng Thế hệ vừa ghi được điểm.

\ct

Đến lượt Korone vào vị trí đánh bóng.

Nàng Thỏ hăng hái đeo găng tay bắt bóng, chỉnh lại tư thế, đứng đúng vị trí chẳng khác nào một tay chơi chuyên nghiệp. Nàng Pháp sư lùi lại một chút để quan sát, còn cậu trọng tài thì bước sang một bên, chắp tay sau lưng như thể đang thưởng thức một màn kịch ngắn đặc sắc.

``Sẵn sàng rồi đó, Pekora-chan\~{} Cẩn thận nha\~{}'' Korone nâng cây gậy lên, ánh mắt tập trung hoàn toàn vào quả bóng đang nằm gọn trong tay đối phương.

*VÚT!*

Pekora tung quả bóng ra với một lực vừa phải. Rushia chỉ vừa kịp hô lên câu ``Đánh nhẹ thôi nha!'' thì\dots

\textbf{*BỐP!!*}

Một cú đánh trời giáng, chính xác từ lần đầu tiên cầm gậy. Quả bóng bay vút ra khỏi khung hình như một viên đạn được bắn ra từ súng.

*XOẢNG!*

Cú đánh đầy uy lực của nàng Khuyển không chỉ phá tung một ô cửa sổ của văn phòng, mà còn kéo quả bóng bay thẳng ra ngoài khuôn viên làm việc.

``\textbf{Peko!? Mạnh thế vậy!?}'' nàng Thỏ hét lên bất ngờ, rồi lập tức phóng người đuổi theo quả bóng.

``\textbf{Chờ tôi với!}'' Rushia cũng gọi theo, chạy ngay sau lưng Pekora. Thấy vậy, Korone cũng vội vàng chạy theo, vừa cười vừa hú lên: ``Khoan khoan, chị đâu có cố ý đánh mạnh đến thế đâu!''

Cả ba cô gái lao ra khỏi khu quay phim như một cơn lốc ba màu: xanh dương, xanh lục và nâu, nhanh đến mức chỉ còn lại vệt bóng thoáng qua.

Kuon đứng im một lát, mắt dõi theo hướng các cô gái vừa chạy đi. Ryuuhou bước thong thả lại gần anh trai, liếc nhìn mảnh kính vỡ dưới sàn rồi bình thản nhận xét: ``Kính cường lực của ô cửa số ba đấy. Onii chuẩn bị sẵn bài giải trình với chị A đi là vừa.''

Cậu Tổng tháo chiếc mũ trọng tài ra, khẽ thở một hơi dài: ``\dots Biết ngay là kiểu gì chuyện này cũng sẽ bất thường mà.''

Cô em gái của cậu nhấp nốt ngụm trà cuối cùng trong ly, khẽ mỉm cười: ``Nhưng mà, ít nhất thước phim vừa quay được cũng đủ chân thực để gửi cho kịch bản của Watame-san rồi đó.''

Cậu Tổng đóng cuốn sổ ghi chép lại, quay gót bước chậm rãi về phía máy quay, như thể đang trò chuyện với cả khán giả (và chính người viết nên câu chuyện này): ``Cảnh kết thúc. Bản dựng này dùng được rồi.''

%==========================================TRUYỆN PHỤ=========================================
\trva

Tập truyện ``Cuộc tập dượt bóng chày'' thực chất được lấy cảm hứng từ các khung hình xuất hiện trong tập phim ``Phi vụ 60 giây'' (chính là {\flaregothic ÉPISODE 127}, tập trước đó.)

Theo các kịch bản mà chính Kanata viết cũng như theo suy đoán của người biên dịch, đây là tập phim có thời lượng đặc biệt ngắn, chỉ kéo dài khoảng 60 giây, được biên dịch lại thành kịch bản bởi Lamy, Kanata và Watame trong một buổi tối\dots\ hỗn loạn, với bối cảnh là chính văn phòng \hll. Truyện sau đó được viết lại vào quyển số Mười - hồi Một bởi người biên dịch và sự giúp sức của ChatGPT. Trong đó, người biên dịch sử dụng ba bức hình của tập 127 để phác họa:

\begin{figure}[H]
    \centering
    \begin{subfigure}[t]{0.45\textwidth}
        \centering
        \includegraphics[width=8cm]{../../../../Image/Book10/c-5a.png}
    \end{subfigure}
    \quad
    \begin{subfigure}[t]{0.45\textwidth}
        \centering
        \includegraphics[width=8cm]{../../../../Image/Book10/c-5b.png}
    \end{subfigure}
    \quad
    \begin{subfigure}[t]{0.45\textwidth}
        \centering
        \includegraphics[width=8cm]{../../../../Image/Book10/c-5c.png}
    \end{subfigure}
\end{figure}

Dù ban đầu đây chỉ là một thử nghiệm nhanh, nhưng câu chuyện ngắn xoay quanh Rushia, Korone và Pekora cùng trò bóng chày bất ngờ lại trở thành một trong những bản dựng có triển vọng nhất, và cũng là nguồn cảm hứng cho bản chuyển thể văn học này.

\end{document}